\documentclass[a4paper,11pt]{article}

\usepackage[utf8]{inputenc}
\usepackage[T1]{fontenc}
\usepackage[german]{babel}
\usepackage[a4paper,margin=2.5cm]{geometry}
\usepackage{enumitem}
\usepackage{titlesec}
\usepackage{longtable}
\usepackage{setspace}
\usepackage{parskip}

\begin{document}

% ---------------------------
% Deckblatt
% ---------------------------
\begin{titlepage}

\centering

\vspace*{3cm}

{\Huge \textbf{Meetingprotokoll}}\\[0.5cm]

{\Large Statusbericht zur Projektarbeit}\\[2cm]

\rule{12cm}{0.5pt}\\[0.5cm]

{\Large Projektarbeit}\\[0.3cm]

{\large 21. August 2026}

\rule{12cm}{0.5pt}

\vfill

{\large THWS / MINOVA}

\end{titlepage}

% ---------------------------
% Inhaltsverzeichnis
% ---------------------------
\tableofcontents
\newpage

% ---------------------------
% Inhalt
% ---------------------------

\section{Allgemeine Informationen}

\begin{tabular}{ll}
\textbf{Datum:} & 21.08.2026 \\
\textbf{Art des Meetings:} & Projektmeeting / Statusbesprechung \\
\textbf{Abgabe Projektdokumentation:} & 16.09.2026 \\
\textbf{Abschlusspräsentation:} & 23.09.2026 um 17:00 Uhr \\
\end{tabular}

\vspace{1cm}

\section{Projektstatus}

Im Rahmen des Projektmeetings wurde die aktuelle Gliederung der
Projektdokumentation vorgestellt und gemeinsam überprüft. Dabei wurde
insbesondere die Struktur der einzelnen Kapitel betrachtet und Feedback
zur inhaltlichen Abgrenzung sowie zur Benennung verschiedener Abschnitte
gegeben.

\vspace{0.5cm}

\section{Besprochene Punkte}

Während des Meetings wurden folgende Punkte zur Gliederung und
Projektdokumentation besprochen:

\begin{itemize}[leftmargin=1.2cm]

\item Das Projektteam soll prüfen, ob in der Einleitung tatsächlich vier
separate Unterpunkte erforderlich sind oder ob eine kompaktere Struktur
sinnvoller ist.

\item Die Kapitel Motivation und Problemstellung sollen
zusammengeführt werden.

\item Kapitel 2 soll ausschließlich als \textbf{IST-Analyse} bezeichnet
werden.

\item Die Darstellung des Prozesses und das zugehörige Diagramm
überschneiden sich inhaltlich. Daher soll geprüft werden, wie diese Inhalte
sinnvoll zusammengeführt werden können.

\item Die Anforderungen an das System sollen in der Dokumentation
klar herausgearbeitet und beschrieben werden.

\item Die Begriffe Frontend und Backend sollen eindeutig
und konsistent verwendet werden, sodass die jeweiligen Systembestandteile
klar voneinander abgegrenzt sind.

\item Kapitel 5 soll als \textbf{Konzeption und Implementierung} strukturiert
werden.

\item Das Frontend soll als separates Kapitel behandelt werden. Innerhalb des
Frontend-Kapitels sollen unter anderem das \textbf{Avisierungsdashboard}
sowie die \textbf{Bestätigungsseite für Kunden} separat beschrieben werden.
Die einzelnen Frontend-Bestandteile sollen als separate Komponenten
umgesetzt und die interne Aufteilung des Frontends in der Dokumentation
erläutert werden.

\item Für den bisherigen Abschnitt zur Fahrerkartenansicht soll eine passendere
und eindeutigere Bezeichnung gewählt werden, beispielsweise
\textbf{Live-Tracking}.

\item Neben dem Ergebnis soll auch beschrieben werden, wie die einzelnen
Komponenten umgesetzt wurden und welches Vorgehen bei der Implementierung
gewählt wurde.

\item \textbf{Projektmanagement} und \textbf{Projektergebnis} sollen als
getrennte Abschnitte behandelt werden.

\item Die bisherigen Punkte 9 und 11 der Gliederung sollen entfernt werden. Die bisherigen Kapitel 8 und 10 sollen zu einem gemeinsamen Kapitel
zusammengeführt werden.

\item Zur besseren Darstellung sollen geeignete Screenshots
in die Dokumentation aufgenommen werden, beispielsweise von versendeten
Nachrichten oder relevanten Benutzeroberflächen.

\item Am Ende der Projektdokumentation sollen die Abschnitte
\textbf{Nutzung von Künstlicher Intelligenz} und
\textbf{Verwendete Hilfsmittel} ergänzt werden.

\item Bei der Abgabe beziehungsweise in der zugehörigen E-Mail soll angegeben
werden, welche Teile der Projektdokumentation von welchen Teammitgliedern
verfasst wurden. Innerhalb der Dokumentation können zur Kennzeichnung der
Zuständigkeiten die Kürzel der Teammitglieder verwendet werden.

\end{itemize}
\vspace{0.5cm}

\section{Nächste Schritte}

\begin{itemize}[leftmargin=1.2cm]

\item Bei Bedarf kann ein weiteres Projektmeeting mit Frau Professor
Saueressig für den 28.08.2026 angefragt werden.

\end{itemize}

\vspace{0.5cm}

\section{Nächste Projektphase}

Die nächste Projektphase konzentriert sich auf die weitere Ausarbeitung
der Projektdokumentation. Dabei soll die Gliederung entsprechend dem
erhaltenen Feedback überarbeitet und die Dokumentation schrittweise
inhaltlich ergänzt und vervollständigt werden. Die weitere Bearbeitung
erfolgt im Hinblick auf die bevorstehende Abschlusspräsentation.

\vspace{0.5cm}

\section{Nächste Meilensteine}

\begin{itemize}[leftmargin=1.2cm]

\item Weitere Ausarbeitung der Projektdokumentation.

\item Überprüfung und Testen der Anwendung sowie Behebung möglicher Fehler.

\item Fertigstellung der Projektdokumentation bis zum 16.09.2026.

\item Vorbereitung der Abschlusspräsentation am 23.09.2026 um 17:00 Uhr.

\end{itemize}

\vspace{1cm}

\section{Zusammenfassung}

Im Projektmeeting am 21.08.2026 wurde die aktuelle Gliederung der
Projektdokumentation gemeinsam betrachtet und besprochen. Dabei wurde
Feedback zur Struktur und weiteren Ausarbeitung der Dokumentation gegeben.

In der nächsten Projektphase liegt der Schwerpunkt auf der Überarbeitung
und Vervollständigung der Projektdokumentation sowie auf der weiteren
Überprüfung der Anwendung. Anschließend erfolgt die Vorbereitung auf die
Abschlusspräsentation.

\end{document}